\documentclass{beamer}
% imprimir
% \documentclass[handout]{beamer} 
% \usepackage{pgfpages}
% \pgfpagesuselayout{4 on 1}[a4paper,landscape,border shrink=5mm]

\mode<presentation> {
  \usetheme{Warsaw}
  \setbeamercovered{transparent}
}

\usebackgroundtemplate{\includegraphics[width=\paperwidth]{format/libresoft-bg.png}}
\usepackage[spanish]{babel}
\usepackage[utf8]{inputenc}
\usepackage{graphics}
\usepackage{amssymb} % Simbolos matematicos

%\definecolor{libresoftgreen}{RGB}{162,190,43}
%\definecolor{libresoftblue}{RGB}{0,98,143}

%\setbeamercolor{titlelike}{bg=libresoftgreen}

%% Metadatos del PDF.
\hypersetup{  
  pdftitle={Presentación del curso},
  pdfauthor={Miguel Vidal},
  pdfcreator={GSyC/Libresoft},
  pdfproducer=PDFLaTeX,
  pdfsubject={Arquitectura de servidores con software libre},
}
%%

\begin{document}

\title{Presentación del curso}
\subtitle{Arquitectura de servidores con software libre}
\institute{\{mvidal,jfcastro\}@gsyc.urjc.es} 
\author{Miguel Vidal, José Castro}
%\date{\today}
\date{25 de marzo de 2011}

\frame{
\maketitle
\begin{center}
\includegraphics[width=6cm]{format/gsyc-urjc}
\end{center}
}

\frame{
~
\vspace{4cm}

\begin{flushright}
{\small
(cc) 2009-2011 Miguel Vidal, José Castro. \\
  Esta presentación se publica bajo una licencia Creative Commons Reconocimiento 3.0 España, disponible en 
  \url{http://creativecommons.org/licenses/by/3.0/es/}
 

\bigskip

}
\end{flushright}
}
%%

\section{Introducción}

\begin{frame}
\frametitle{Información general}
\begin{itemize}
\item \alert{Objetivo:} ofrecer una panorámica de las soluciones que ofrece el software libre para construir complejas arquitecturas virtualizadas.
\item \alert{Inicio:} 25 de marzo 2011 -- Fin: 10 de junio 2011
\item \alert{Horario:} los viernes de 16:30h a 20:30h 
\item \alert{Lugar:} Madrid On Rails
\item \alert{Material:} ordenador y libro de referencia
\item \alert{Título:} título propio de la URJC. Al finalizar, obtendréis el certificado de aprovechamiento del curso.
\end{itemize}
\end{frame}

\section{Temario}
\begin{frame}
\frametitle{Clases}
\begin{enumerate}
\item 25 de marzo: Introduccion
\item 1 de abril: Seguridad
\item 8 de abril: SOs libres, homogeneización
\item 29 de abril: Redes
\item 6 de mayo: Storage as a Service (StaaS)
\item Servicios de Internet
\item Virtualización I
\item Virtualización II
\item Cloud Computing
\item Clusters de Alta Disponibilidad
\end{enumerate}
\end{frame}


\begin{frame}
\frametitle{Conferencias invitadas}

\begin{enumerate}
\item Introducción al software libre (Libresoft)
\item IPv6 (Eva Castro, GSyC)
\item Monitorización con Pandora (Ártica)
\item Certificación con CaCERT (a cargo de un notario)
\item La experiencia con software libre de empresas TIC: Andago y OpenSistemas
\item Servidor web Cherokee (a cargo del líder del proyecto)
\item Evolución del CPD, convergencia de almacenamiento y virtualización (Víctor Fernández, SIA Group)
\item Virtualización en SPARC con LDOMs (Gnome Hispano)
\item Tutorial de OpenNebula (a cargo del OpenNebula core Team)
\item Aspectos legales del cloud computing (departamento jurídico de Google)
\end{enumerate}

\end{frame}


\section{Metodología}
\begin{frame}
\frametitle{Clases}

\begin{itemize}
\item Clases. 
\item Libro de referencia
\item En casa: 
\item En clase: 
\item Foros en Moodle
\end{itemize}

\begin{center}
\url{Moodle}{http://moodle.libresoft.es}
\end{center}

\end{frame}


\begin{frame}
\frametitle{Equipos y material}

\begin{itemize}
\item Puedes traer tu propio portátil
\item WIFI para el entorno:
\begin{itemize}
\item ESSID:   
\item usuario: 
\item contraseña: 
\end{itemize}
\item HashTag Twitter: \#casul (Curso Arq. Sistemas Urjc Libresoft)
\item Canal IRC: \#casul en irc.freenode.net
\end{itemize}

\end{frame}

\begin{frame}
\frametitle{Contacto}
\begin{itemize}
\item E-mail: \{mvidal,jfcastro\}\@gsyc.urjc.es
\item Twitter: @mvidallopez y @jfcastroluis
\item Moodle: \url{http://moodle.libresoft.es}

\end{itemize}
\end{frame}

\begin {frame}
\frametitle{Preguntas}
\begin{center}
\huge{Preguntas y dudas}
\end{center}
\end{frame}

\begin {frame}
\frametitle{Preguntas}
\begin{center}
\huge{Os toca presentaros}
\end{center}
\end{frame}


\end{document}
